\chapter[Sermon 82. Of the Marriage of the Lamb, and of the Making Ready of the Lamb's Wife.]{Of the Marriage of the Lamb, and of the Making Ready of the Lamb's Wife.}
\label{sermon:082}
\markright{Sermon 82. Of the Marriage of the Lamb, and of the Making Ready ...}

And his wife prepared herself. And it was decreed that she should be clothed with pure and beautiful silk, for silk is the righteousness of the saints.

The saints celebrate the Lord with praises, rejoicing, and hymns. There are countless reasons for this, yet two stand out above the others. The first is that the Lord has judged the harlot and avenged the blood of the saints. The second is that the marriage of the Lamb has come. They rejoice therefore at the justice of God, by which he has punished the ungodly, and at his mercy or grace, by which he gives to the godly a blessed life. But here we must speak of the marriage.

\index{Turks}There is much mention made in the holy Scriptures, both of the Old and New Testaments, of matrimony and marriage. This cannot be explained according to the literal meaning, but by an allegory, lest we fall into shameful and most outrageous absurdities alongside the Turks and Mahometans. For spiritual things are represented by physical matters. The spiritual essence of this is: God the Father, the lover of mankind, will save men by his Son. This is declared by a parable of wedlock and marriage. And in matrimony there is a contract or a binding agreement, there is coupling or a handfasting of each party, and finally, marriage.

\index{Ezekiel (Book of)}In the contract not only the young man and the maid are affianced, but also the whole manner of the marriage to come is appointed, and an order taken. For the lawyers say that affiancing is a promise of the marriage to come. This contract was made at the beginning of the world, where God promises that he will deliver mankind by his son, and receive him into glory. Hereunto appertain all the promises of Christ, of the remission of sins, and everlasting life. Moreover, the duties of the spouse are prescribed. She promises to be obedient, and other things, \&c. Christ, the son of God the Father, as bridegroom, affiances himself to all the chosen through his free grace: he promises them his righteousness, all heavenly gifts, and eternal life. He takes upon himself moreover all the infirmities of the bride, and purges her filthiness. And the bride is affianced to him by faith, as it is with Hosea, and binds herself wholly to him: after whose will and law she frames herself wholly. For she is the body of a lively head. As Paul says in 5.1 to the Ephesians. The bride’s leaders are the prophets, patriarchs, Apostles. So John the Baptist in 3.29 of John, calls himself the friend of the bridegroom. He adds, to be the spouse of Christ. Paul, 2 Corinthians 11, “I have espoused you to one husband, a chaste virgin,” \&c. Hereunto the 16th chapter of Ezekiel seems to appertain.

The joining together of either party is performed after they are betrothed, with certain ceremonies: namely, by taking each other by the hands, and certain words spoken, a token or a ring is given, and so forth. But immediately after the beginning, a league or bond was made between God and men, which is often read of, not without ceremonies, certain words and sacrifices being offered, as by Abraham, Moses, and others. God binds himself to men, and men to him, and that not without sacraments. Hereunto belong all those things that God would be in league with man, and have men bounden to him, and all his things communicated to us. And this marriage, of all others, is most strictly joined while the Son of God has united our flesh into one and the same person with him, and has commanded the Apostles to preach unto all that he will have a communion with the faithful. Of this communion many things are read everywhere in the scriptures. And he has given a pledge of faith and perpetual amity, not a ring of gold, but rather the sacraments: yea, even the Holy Spirit, as Paul says in 1 Corinthians 1 and to the Ephesians 1.

\index{Resurrection}And the marriage shall be celebrated in the resurrection of the dead. The souls truly pass from bodily death into everlasting life; but the full restoration and salvation of humankind is not made perfect unless the body also comes. Therefore, at the resurrection comes the marriage of the Lamb, that is, of Christ our Redeemer. Then we are carried to meet Christ in the air, and he brings his wife into the bridal chamber of eternal glory and blessing; then shall be held that feast and delightful supper, then shall the bride enjoy forever the love of the bridegroom. This shall truly be the marriage of the Lamb. And the marriage shall be the merrier, for that the harlot being cast out and condemned, the wife and virtuous matron shall have the full and perfect joy alone. At this joy and at this marriage, the holy inhabitants of heaven rejoice.

\index{Antichrist}Moreover, the saints experience here also a certain preparation of the Bride, so that by the way the godly may understand what is best for them and to what they should dedicate themselves in the last age. Let us prepare ourselves to meet the Bridegroom. For we look for the judge every hour. And we prepare ourselves not in one hour or day, but throughout our lifetime. And how we should be prepared, the Lord himself shows by the parable of the ten virgins. Let us adorn ourselves with true faith against Antichrist in the later days. Let us beautify ourselves with the works of charity, the works also of righteousness, chastity, and temperance. Let us not be corrupted and defiled with drunkenness, bloodshed, and the cares of this world.

\index{Justification}Furthermore, lest anyone should ascribe this preparation to his own merit, strength, and virtue, and that we should see also that this preparation chiefly consists in providing the garment, John adds immediately, “And to her was granted that she should array herself.” If it is given, then it is not prepared by our force or means, 1 Corinthians 4. If it is given, then it is not bought by popish dealings. Read Acts 8. He also expresses the kind of garment, of clean or pure silk, and shining or bright. For we read also in the gospel of the wedding garment. The Apostle frequently exhorts us that we should put on the Lord Jesus. These things are in allegory. But he immediately explaining this kind of garment, says that silk is the righteousness of saints. Saints he calls the faithful. But where there is only one justification through faith in Christ, John speaks of justifications in the plural number. For those who are freely justified through faith alone, immediately perform sundry and many works of righteousness. For he who is righteous, as John says, does works righteousness. Therefore there are justifications, namely, the righteousness of faith justifying, and the righteousness of works justifying: that is to say, declaring us to be justified by faith alone. For we are purified by the blood of Christ freely, which we receive by faith, and are fully justified, as Paul witnesses in Romans 3. Again, those who are righteous do sundry works of righteousness and commend themselves to God. Thus they do not appear naked, but clothed with their wedding garment, as we touched upon in the third chapter of this book.

\index{Daniel (Book of)}And it is very fitting that the bride’s garment is called pure or clean, not because of her own person, whom we know to be always hindered and weakened by the flesh, but because of the Spirit sanctifying and the blood of the Son of God. As Paul testifies in Ephesians 5:1, and in 1 John 1:1. The garment is also said to be shining and bright, and that for the glorifying of the saints to come. This is mentioned in Daniel 12 and Matthew 13. For righteousness leads to glory. For whom he has justified, he has also glorified. To him be praise, honor, and glory.
